% =====================================================================
% FULL REPORT DRAFT — Stage 1 complete, Stage 2 placeholders marked TODO
% Requires: \usepackage{amsmath}. Tables use plain \hline (no booktabs).
% Formatted for the UZH thesis template: \chapter top-level, \section next,
% \subsection below that. Cross-references use \autoref. Acronyms are not
% yet wrapped in \gls{}/\glspl{} — that requires the actual keys defined in
% your resources/glossary.tex, which this session does not have access to.
% =====================================================================


% =====================================================================
% 0-abstract-en.tex
% =====================================================================
\chapter*{Abstract}

We compare three strategies for dialogue-act classification of Reddit
ChangeMyView (CMV) parent-reply pairs into four classes (\texttt{agree},
\texttt{disagree}, \texttt{question}, \texttt{statement}): (A) direct
annotation with a decoder-only LLM (Llama-3.1-8B-Instruct, few-shot
prompting), (B) a RoBERTa-base encoder fine-tuned on LLM-generated silver
labels, and (C) a RoBERTa-base encoder fine-tuned directly on 300
human-labeled gold parent-reply pairs. All three are tuned on the same
gold set and evaluated once on an independent, expanded 332-pair balanced
held-out set.
Strategy A and B are statistically indistinguishable on held-out
(macro-F1 0.673 vs.\ 0.669, paired-bootstrap effect size $\Delta=0.004$,
practically a tie rather than an underpowered comparison), while Strategy C
trails both (macro-F1 0.586) with seed-to-seed variance an order of
magnitude larger than originally estimated from three seeds. Across all
three strategies, \texttt{question} is consistently the easiest class and
\texttt{statement} the hardest. A detailed error analysis finds a
bidirectional \texttt{agree}/\texttt{disagree} confusion specific to the two
RoBERTa-based strategies (B and C) and absent in the LLM-based strategy (A),
suggesting a backbone-specific limitation rather than inherently ambiguous
text. This report covers Stage 1 in full. Stage 2 is in progress, and its
plan is outlined in the Future Work chapter.


% =====================================================================
% 1-introduction.tex
% =====================================================================
\chapter{Introduction}
\label{k:einleitung}

Understanding the function of a reply in an online argumentative discussion
--- whether it agrees, disagrees, asks a question, or simply states a claim
--- is a prerequisite for downstream analysis of persuasion and disagreement
dynamics. Reddit's r/ChangeMyView (CMV) is a natural testbed for this:
parent-reply pairs are argumentative by design, and the four-class scheme
used here (\texttt{agree}, \texttt{disagree}, \texttt{question},
\texttt{statement}) captures the coarse dialogue act of a reply relative to
its parent.

Three practical routes exist for building a classifier of this kind.
(A) Prompt a large pretrained LLM directly, with no task-specific training
data. (B) Use that LLM to label a larger pool of unlabeled data, then
fine-tune a small, fast encoder on those \emph{silver} labels. (C) Skip
the LLM entirely and fine-tune the small encoder directly on a limited
pool of \emph{gold} (human-labeled) data. Each route makes a different
trade-off
between annotation cost, inference cost, and expected accuracy, and it is
not obvious in advance which wins for a task with only a few hundred gold
pairs available.

This report addresses three research questions: (RQ1, data efficiency) is
300 gold pairs enough for direct fine-tuning to match LLM annotation? (RQ2,
distillation) does distilling an LLM's own labels into a fine-tuned model
improve on using the LLM directly? (RQ3, label source) given a fixed
encoder backbone, does fine-tuning on gold labels or LLM-generated silver
labels perform better? All three strategies are tuned on the same
300-pair gold set and evaluated once, after locking each strategy's final
configuration, on an independent held-out set --- expanded during this
project from an initial 80-pair pilot to a 332-pair balanced set
specifically to gain enough statistical power to address these questions.

This project proceeds in two stages. \textbf{Stage 1}, reported here in
full, establishes and evaluates the three strategies above. \textbf{Stage 2}
is in progress. Its plan is described in the Future Work chapter and will
be integrated into the Methodology and Experiments chapters once complete.


% =====================================================================
% 3-related-work.tex
% =====================================================================
\chapter{Related Work}
\label{k:relatedwork}

% TODO: this chapter is a topic skeleton, not a literature review — no
% citations have been verified in this session and none should be
% fabricated. Please fill in citations for each area below, using
% \textcite{} / \parencite{}.

\section{Dialogue Act / Speech Act Classification}
% TODO: cite prior dialogue-act tagging work (e.g., Switchboard-style
% schemes, argument-mining dialogue-act taxonomies) and situate the
% agree/disagree/question/statement scheme used here relative to them.

\section{LLM Zero-/Few-Shot Annotation}
% TODO: cite work on using LLMs as annotators/labelers, prompt-based
% classification without task-specific fine-tuning, and known failure
% modes of few-shot prompting for classification tasks.

\section{Knowledge Distillation via Silver Labels}
% TODO: cite work on training smaller models on LLM-generated labels
% (silver-label distillation), and any prior findings on when a distilled
% student matches or fails to match its LLM teacher.

\section{Fine-Tuning Encoders on Small Labeled Sets}
% TODO: cite RoBERTa/BERT fine-tuning literature, especially work on
% training instability and warmup/epoch sensitivity with small (~hundreds
% of instances) datasets, relevant to the instability found in Strategy C
% (\autoref{sec:strategyc}).

\section{Argument Mining on r/ChangeMyView}
% TODO: cite prior CMV-specific work (e.g., the Winning Arguments Corpus
% this project's data is drawn from) and any existing dialogue-act or
% persuasion-related annotation schemes used on this corpus.


% =====================================================================
% 4-methodology.tex
% =====================================================================
\chapter{Methodology}
\label{k:methodology}

\section{Stage 1 Methodology}

\subsection{Task and Label Scheme}

Stage 1 classifies the dialogue act of each reply in a CMV (Winning
Arguments Corpus) parent-reply pair, using a four-label scheme:
\texttt{agree}, \texttt{disagree}, \texttt{question}, and
\texttt{statement}. These labels capture the surface communicative
function of a reply relative to its parent. We compare three modelling
strategies:

\begin{itemize}
\item \textbf{Direct LLM Annotation} (Strategy A): LLM-predicted
      dialogue-act labels are used directly, via zero-/few-shot
      prompting, with no fine-tuning. This establishes whether prompting
      alone is sufficient for this classification.
\item \textbf{Silver-Standard Fine-Tuning} (Strategy B): the LLM is
      used to label a larger pool of unlabeled CMV data, and the
      resulting silver-standard annotations are used to fine-tune a
      RoBERTa-base classifier. This tests whether distilling the LLM's
      own annotations into a smaller model matches or exceeds using the
      LLM directly, at lower inference cost.
\item \textbf{Gold-Standard Fine-Tuning} (Strategy C): RoBERTa is
      fine-tuned directly on the 300-pair gold-standard annotations,
      using cross-validation. This provides a data-efficiency reference:
      whether 300 human-labeled pairs alone are enough to train a usable
      classifier, without an LLM in the loop.
\end{itemize}

All three strategies are tuned on the gold-standard set and evaluated
once, after locking each configuration, on an independent held-out set
(\autoref{sec:data}). Model choices, hyperparameters, and results for
each strategy are reported in \autoref{sec:experiments}.

\subsection{Gold Set (Dev) and Held-Out Test Set}
\label{sec:data}

The gold set (\texttt{dev}) consists of 300 parent-reply pairs: 145
disagree, 64 agree, 47 statement, and 44 question. It is used exclusively
for tuning, and all three strategies select their final configuration on
it. The pairs were double-annotated by two annotators and reconciled by
discussion (\autoref{sec:iaa}).

The held-out test set (\texttt{TEST\_CSV}) is touched only once per
strategy, after configuration lock. It was expanded from an 80-pair pilot
(\textbf{pilot80}) to a 332-pair, 83-pairs-per-class balanced set for
statistical power (\autoref{sec:no-ranking}): candidates were sampled
from the Winning Arguments Corpus excluding gold/silver overlaps, labeled
into a 530-pair pool, then downsampled to the smallest class
(\texttt{statement}).

No identical (Parent, Reply) pair appears across gold, silver, and
held-out (verified programmatically). Different replies to the same
parent/thread are intentionally permitted, reflecting real Reddit
structure. An audit found 262 of 276 traceable gold pairs share a thread
with the silver pool, so Strategy B results reflect in-domain, same-topic
performance rather than unseen-topic performance.

\subsection{Inter-Annotator Agreement}
\label{sec:iaa}

Both gold-300 and held-out were double-annotated by two annotators and
reconciled by discussion. In both cases, the pre-reconciliation labels
were not retained, so only raw agreement can be reported. Cohen's kappa is
not computable. Raw agreement was 89.67\% on gold-300 (300 pairs) and
95.18\% on held-out (332 pairs, 16 disagreements).

\section{Stage 2 Methodology}

% TODO: to be written once Stage 2's design is finalized. Expected to
% extend the protocol above (see Future Work, \autoref{sec:future-work})
% rather than replace it.


% =====================================================================
% 5-experiments.tex
% =====================================================================
\chapter{Experiments}
\label{sec:experiments}

\section{Stage 1 Experiments}

All results in this section are Stage 1, for the three strategies
introduced in \autoref{k:methodology}. Scores are reported side by side
but should not be read as a ranking of ``best'' strategy --- see
\autoref{sec:no-ranking}.

\subsection{Strategy A: LLM Annotation}

% TODO: touvron2023llama is the Llama 2 paper (arXiv 2307.09288).
% Replace with the correct Llama 3.1 technical report citation once
% available. This is a placeholder, not a verified match.
Strategy A uses Llama-3.1-8B-Instruct \parencite{touvron2023llama} at
temperature 0, comparing zero-shot against few-shot prompting. Few-shot
clearly wins on the gold set and is locked as the final configuration,
reaching macro-F1 0.673 on held-out (\autoref{table:strategya-results}).

The dominant error, on both dev and held-out, is that \texttt{disagree}
pairs are frequently misclassified as \texttt{statement}, while
\texttt{question} is comparatively easy (per-class breakdown in
\autoref{table:strategya-results} and \autoref{table:strategya-heldout-cm}).
Notably, \texttt{agree} and \texttt{disagree} are rarely confused with
each other on held-out --- both instead leak into \texttt{statement}
--- a pattern that differs qualitatively from Strategies B and C
(\autoref{sec:error-analysis}).

\subsection{Strategy B: RoBERTa on Silver Labels}

Strategy B fine-tunes RoBERTa-base on silver labels generated by
Strategy A's LLM. Two settings control this: \texttt{sample\_seed} picks
which pairs are drawn from the silver pool, and \texttt{train\_seed}
controls training randomness. The train/validation split is grouped by
\texttt{source\_root}, so no Reddit thread appears on both sides.

To choose how much silver data to use, we tested pool sizes from 500 to
10{,}000 pairs (full curve in \autoref{app:extra-matrices}). Macro-F1
rises sharply up to 1000 pairs, peaks at 2500 pairs, then declines slowly
--- more silver data does not keep helping. We locked the configuration
at 2500 pairs, \texttt{sample\_seed} 123, with class weights enabled.
Class weights give a small, consistent improvement over no weights
(0.6179 vs.\ 0.6113 macro-F1), though this difference has not been
tested for significance. Full results are in
\autoref{table:strategyb-results}, and the dev confusion matrix is in
\autoref{table:strategyb-dev-cm}.

On held-out, macro-F1 is 0.6686 (kappa 0.5649), averaged across three
training seeds. Its weakest class is now \texttt{statement}, not
\texttt{disagree} as on the earlier, smaller held-out set --- so B's
weakest class is not stable across evaluation sets (full confusion
matrix in \autoref{table:strategyb-heldout-cm}).

Because most gold pairs and the silver pool come from the same Reddit
threads (\autoref{sec:data}), these results reflect in-domain
performance, not performance on genuinely unseen topics.

\subsection{Strategy C: RoBERTa on Gold Labels}
\label{sec:strategyc}

Strategy C uses 5-fold cross-validation on the 300 gold pairs: the data
is split into 5 parts, and each part is predicted once by a model
trained on the other 4. Each seed's 5 predicted parts are combined into
one 300-pair \emph{out-of-fold (OOF)} set, and macro-F1 is computed once
on that combined set, not averaged across the 5 parts separately.

Early attempts failed: the model would ignore one or more classes
entirely (\emph{class collapse}), giving macro-F1 of only 0.16--0.20. To
diagnose this, we checked whether each model could even fit its own
training data --- mostly, it could not (train-fit macro-F1 0.326,
strongly correlated with eval quality). This points to unstable
training, not too little data. Adding a warmup period before ramping up
the learning rate (20 steps), and training longer (8 epochs instead of
fewer), fixed this: with class weights enabled, no run collapsed a
class, and OOF macro-F1 reached 0.5309 (3 seeds). Without class weights,
training was still unstable (8 of 15 runs collapsed a class).

A full grid search over 9 warmup/epoch combinations confirmed
warmup=20, epochs=8 as the best setting, and it was locked as final.
Training for only 6 epochs is undertrained at every warmup value tested.
10 epochs scores about as well on average but is 8--13$\times$ noisier
across seeds, so less reproducible.

To check whether the agree/disagree confusion in
\autoref{sec:error-analysis} was just bad luck with one seed, we ran 5
more seeds on the same locked configuration
(\autoref{table:strategyc-variance}). Across all 8 seeds, macro-F1
averages 0.543 with a standard deviation of 0.037 --- about ten times
larger than the original 3-seed estimate. The original three seeds
happened to land unusually close together by chance, and the true
variance is much bigger. The warmup/epoch choice is unaffected (still
the best mean, still no class collapse), but the earlier figure,
``0.5309 $\pm$ 0.0038,'' should not be read as evidence that training is
highly stable. Dev-vs-held-out results are compared in
\autoref{table:strategyc-dev-vs-heldout}.

On held-out, the original three seeds average macro-F1 0.5859, but
spread widely: 0.4874 (95\% CI [0.43, 0.54]), 0.6269 ([0.58, 0.67]), and
0.6435 ([0.59, 0.69]). The worst seed here is the same seed that scored
worst on dev, which fits the 8-seed variance above --- not a one-off
anomaly (full held-out confusion matrices per seed in
\autoref{app:extra-matrices}).

\subsection{Cross-Strategy Comparison}
\label{sec:no-ranking}

The ranking is not stable across evaluation sets (\autoref{table:cross-strategy-comparison}).
On gold, B $>$ A $>$ C. On the earlier 80-pair held-out set, A $>$ C $>$ B.
On the new 332-pair held-out set, A $\approx$ B $>$ C, with A's and B's
confidence intervals overlapping almost completely.

\texttt{question} is the easiest class and \texttt{statement} the hardest for
all three strategies, consistently across both held-out sets
(\autoref{table:per-class-f1}).

To check whether these differences are likely real, rather than just
random noise from which 332 pairs happened to be drawn, we use a test
called a \emph{paired bootstrap}: the 332 held-out pairs are repeatedly
re-sampled at random, and each strategy's macro-F1 is recalculated on
every re-sample. Comparing two strategies on the same re-sampled pairs
(rather than as if they were two unrelated samples) makes the test more
sensitive, since it cancels out the effect of some pairs simply being
harder than others.

\autoref{table:power-analysis} reports, for each pair of strategies: the
observed difference in macro-F1 ($\Delta$), a 95\% confidence interval
for that difference (the range within which the true difference most
likely falls), whether the difference is statistically significant, the
test's current \emph{power} (roughly, the probability that a test this
size would detect a real difference of this magnitude if one exists),
and how many labeled pairs per class would be needed to reach the
conventional 80\% power threshold (that last column is an analytical
projection from the current sample, not a new experiment).

None of the three comparisons is statistically significant yet, but for
two different reasons. A and B differ by almost nothing ($\Delta =
0.0041$): reaching 80\% power for this comparison would require about
34,900 labeled pairs per class, which is not realistic for this project.
The honest reading is that A and B perform about the same, not that more
data would eventually reveal a difference. A-vs-C and B-vs-C, in
contrast, show a real, if modest, gap ($\Delta \approx 0.04$--$0.05$)
that would become detectable at roughly 330--350 labeled pairs per
class --- well above an earlier estimate of 60--70, but still within
reach of further annotation. \autoref{c:conclusions} discusses why these
results are not condensed into a single ranking.

\subsection{Error Analysis: Confusion Overview}
\label{sec:error-overview}

Pooling errors across all three held-out seeds for each strategy (996
predictions per strategy), the single largest error types are not
\texttt{agree}/\texttt{disagree} confusions, but two other pairs:
\texttt{disagree}$\to$\texttt{statement} (Strategy B's largest error
type, 62 instances) and \texttt{statement}$\to$\texttt{agree} (Strategy
C's largest, 84 instances; also present in B, 56 instances). Reading a
sample of each, together with simple text statistics (word count, and
the rate of hedge words like ``but''/``however'' and negation words like
``not''/``never''), suggests three distinct, specific patterns.

\textbf{\texttt{disagree}$\to$\texttt{statement} (Strategy B).} These
replies disagree by stating a counter-fact or counter-hypothetical, not
by using disagreement language: hedge-word rate is 0.03 (vs.\ 0.13 for
\texttt{disagree} generally) and negation rate is 0.31 (vs.\ 0.51).
For example: \emph{``There is no such thing as cultural decay, only
change.''}; \emph{``Consider this: if this effect did not occur, we
would be heading towards a Malthusian catastrophe.''} The disagreement
is carried entirely by the content, not by any lexical marker, so the
model defaults to reading these as neutral statements.

\textbf{\texttt{statement}$\to$\texttt{agree} (Strategies B and C).} 36\%
of these misclassified replies are five words or fewer, far above
baseline. These are short, reactive replies --- thanks, corrections,
brief asides --- that carry a positive or polite tone without
substantively engaging the argument. For example: \emph{``Whoops,
thanks!''}; \emph{``changed.''}; \emph{``Ahh. Yes I misinterpreted your
comment. Thanks for the clarification.''} The model appears to read
positive or polite tone in a short reply as \texttt{agree}, regardless
of whether it substantively agrees with anything.

\textbf{\texttt{statement}$\to$\texttt{disagree} (Strategy C, 75
instances).} These replies are longer than baseline (15.1 vs.\ 12.7
words) with elevated hedge-word (0.23 vs.\ 0.17) and negation (0.31 vs.\
0.25) rates --- they read as skeptical or challenging in tone (embedded
rhetorical questions, ``but'' qualifiers) without being a genuine
disagreement with the parent's core claim.

Together with the \texttt{agree}/\texttt{disagree} pattern below, this
points to a broader theme: all of B's and C's largest error types
involve replies where the dialogue act is carried by tone, structure, or
implication rather than by an explicit lexical marker, and RoBERTa
fine-tuned on 300--2500 pairs struggles specifically with this class of
example.

\subsection{Error Analysis: Additional Text-Feature Checks}
\label{sec:error-textfeatures}

To complete the negation/contractions/length/spelling check, a
contraction rate and a spelling-error rate (nltk words-corpus heuristic:
a token not found in the corpus and not capitalized) were computed the
same way as the negation and hedge-word rates above, comparing each
error pair against its class baseline.

\textbf{Contractions and spelling: no effect.} Neither rate differs
consistently between misclassified and baseline groups, across any of
the confusion pairs checked
(\texttt{disagree}$\to$\texttt{statement}, \texttt{statement}
$\to$\texttt{agree}, \texttt{statement}$\to$\texttt{disagree},
\texttt{agree}$\leftrightarrow$\texttt{disagree}) or across strategies
overall. This rules out contractions and spelling as drivers of these
errors.

\textbf{Very long replies ($\geq$30 words): elevated specifically for
Strategy A.} 7.3\% of A's misclassified held-out pairs are $\geq$30
words, vs.\ 0.9\% of its correctly classified pairs (roughly 8$\times$)
--- a pattern not present for B (0.9\% vs.\ 4.0\%, inverted) and only
mild for C's \texttt{statement}$\to$\texttt{disagree} pair (6.7\% vs.\
2.4\% baseline, which confirms the ``longer than baseline'' description
above quantitatively).

\textbf{Very short replies ($\leq$5 words): confirms the
\texttt{statement}$\to$\texttt{agree} finding above.} 35.7\% (B) and
32.1\% (C) of \texttt{statement}$\to$\texttt{agree} misclassifications
are $\leq$5 words, vs.\ 25.3\% baseline --- consistent with the
$\sim$36\% figure already reported. Not elevated for Strategy A overall
(12.7\% vs.\ 18.9\% for correctly classified pairs).

\textbf{Ambiguous wording} remains a qualitative judgment. No reliable
automatic proxy was found, so this criterion is assessed only through
the manually read examples above, not a quantitative metric.

This completes the full negation/contractions/length/spelling/ambiguous-
wording checklist: negation and short-reply length show real effects
(above), contractions and spelling do not, very-long-reply length shows
an effect specific to Strategy A, and ambiguous wording is assessed
qualitatively only.

\subsection{Error Analysis: Agree/Disagree Confusion}
\label{sec:error-analysis}

Beyond the confusion pairs above, a smaller but more \emph{reproducible}
problem stands out: a bidirectional
\texttt{agree}/\texttt{disagree} confusion, present in all 8 of Strategy
C's seeds without exception and, to a lesser degree, in Strategy B.
Severity varies by seed (Strategy C's \texttt{agree} recall ranges
0.36--0.59), but direction is constant: \texttt{agree} misclassifies
mainly as \texttt{disagree} in every seed, and vice versa.

This rules out seed-ensembling as a fix: ensembling averages out errors
that are independent across members, but here every seed errs in the same
direction, so a majority vote would preserve the shared bias rather than
cancel it.

To test whether this reflects the text itself or the RoBERTa training
approach, the 27 gold pairs misclassified by all or most of Strategy C's 8
seeds were checked against Strategy A and B's predictions on the same
text (\autoref{table:hard-pairs-accuracy}).

Strategy A performs on these pairs as well as on its overall baseline
(51.9\% vs.\ 49.8\%) --- no elevated difficulty. Strategy B performs 16
points worse than its own baseline (39.5\% vs.\ 55.7\%), and Strategy C
fails on all of them by construction. Since only the RoBERTa-based
strategies are affected, this points to a limitation of RoBERTa
fine-tuning itself rather than inherently ambiguous text --- unlabelable
text should have hurt Strategy A too.

A closer reading suggests a candidate explanation: many pairs show a
concede-then-pivot structure, e.g.\ ``I understand that I'm simplifying...
\emph{But} my view is...'' or ``Fair enough... \emph{Though} I
believe...'', where the true stance is in the second clause. RoBERTa-base,
fine-tuned on only 300--2500 pairs, may lack signal to weight the pivot
over the opening concession. Errors persist even with explicit markers
(``I agree'', ``Indeed''), suggesting an attention-allocation failure
rather than pure ambiguity.

Two alternatives were checked. Sarcasm, hostility, and contempt
annotations show no meaningful difference between the 27 hard pairs and
the overall population ($\sim$14--15\% sarcasm both), ruling this out.
Class imbalance (145 \texttt{disagree} vs.\ 64 \texttt{agree}, 2.3:1) has
not been ruled out, and may independently explain lower \texttt{agree}
recall.

This is a model-cross-validation proxy, not a direct test of these 27
pairs. Gold-300 has measured overall IAA (89.67\%, \autoref{sec:iaa}),
but pre-reconciliation labels were not retained, so these specific pairs
cannot be checked against known disagreements. A blind re-labeling
(Future Work, \autoref{sec:future-work}) is the direct test.

\section{Stage 2 Experiments}

% TODO: to be written once Stage 2 runs are complete.


% =====================================================================
% 8-conclusion.tex
% =====================================================================
\chapter{Conclusion}
\label{c:conclusions}

Stage 1 compared three strategies for four-class dialogue-act
classification of CMV reply pairs on an expanded, balanced, independent
held-out set, but the comparison is not yet complete: of the three
pairwise comparisons (\autoref{sec:no-ranking}), only one is
resolved.

\textbf{A vs.\ B} (does distillation pay off?) is resolved: despite
receiving by far the largest tuning budget of the three strategies, B
ties rather than loses to A (macro-F1 0.669 vs.\ 0.673). The effect size
is near zero, so this is a genuine tie, not a result waiting on more
data --- silver distillation does not detectably beat direct LLM
annotation.

\textbf{A vs.\ C} (data efficiency) and \textbf{B vs.\ C} (silver vs.\
gold source) remain open. Both show a real effect (macro-F1 gaps of
0.087 and 0.083, respectively) but are not yet statistically significant
at the current held-out size (83 pairs/class). Closing them would
require annotating toward $\sim$330--350 confirmed pairs/class
(\autoref{sec:no-ranking}), a substantial effort that is not planned
within the current timeline. Both comparisons are therefore expected to
remain open in this report. The compute-only follow-ups in Future Work
(additional seeds, backbone swaps) can still sharpen the picture without
new annotation.

Two further findings qualify how these numbers should be read. Strategy
C's seed-to-seed variance was understated by roughly $10\times$ when
estimated from only three seeds, a caution for how future small-sample
variance estimates in this project, including Stage 2, should be
interpreted. And a systematic bidirectional \texttt{agree}/\texttt{disagree}
confusion appears in both RoBERTa-based strategies but not in A, evidence
that model backbone --- not just label source --- drives part of what
looks like a strategy-level effect (\autoref{c:limitations}).


% =====================================================================
% 7-future-work.tex
% =====================================================================
\chapter{Future Work}
\label{sec:future-work}

\section{Stage 1 Follow-Ups}

Compute-only follow-ups (no new annotation required):

\begin{itemize}
\item Add more seeds/runs to Strategies A and B to check for previously
      underestimated variance, mirroring what was found for Strategy C.
\item Test the class-imbalance hypothesis: oversample \texttt{agree} and
      rerun Strategy C, check whether the confusion eases.
\item Swap \texttt{roberta-base} for \texttt{roberta-large} in Strategies B
      and C (no new grid search), and/or evaluate a second LLM for
      Strategy A, to test whether the A-vs-B/A-vs-C comparison is
      sensitive to backbone choice.
\end{itemize}

Annotation-dependent follow-ups (not currently planned, given time
constraints):

\begin{itemize}
\item Blind re-label the 27 agree/disagree ``hard'' pairs
      (\autoref{sec:error-analysis}) to test whether real annotation
      disagreement exists, or whether the confusion is purely a model
      limitation --- small-scale (27 pairs) relative to the item below.
\item Annotate toward the $\sim$330--350 confirmed pairs/class identified by
      the power analysis (\autoref{sec:no-ranking}) as the threshold
      needed to resolve A-vs-C and B-vs-C at 80\% power. Note this
      threshold is not expected to resolve A-vs-B, whose effect size is
      near zero. This is the only way to make A-vs-C and B-vs-C
      conclusive. Without it, both remain open questions in this report.
\end{itemize}

\section{Stage 2}

Stage 2 is not yet complete. Its plan will be detailed here once finalized.


% =====================================================================
% 6-limitations.tex
% =====================================================================
\chapter{Limitations}
\label{c:limitations}

\begin{itemize}
\item \textbf{Gold/silver domain overlap (Strategy B).} See
      \autoref{sec:data} --- held-out results reflect in-domain
      (same-thread) performance, not fully unseen-topic performance.
\item \textbf{Gold-300's IAA is aggregate-only.} 89.67\% raw agreement is
      known (\autoref{sec:iaa}), but pre-reconciliation labels were not
      retained, so pair-level annotator disagreements cannot be recovered
      --- the agree/disagree ``ambiguity'' analysis
      (\autoref{sec:error-analysis}) remains a model-cross-validation
      proxy, not a rigorous per-pair IAA test.
\item \textbf{Model family is bound to strategy.} Strategy A uses a
      decoder-only LLM. B and C fine-tune the same encoder (RoBERTa-base)
      on different label sources. A gap between A and \{B, C\} is
      therefore consistent with the LLM being better suited to the task,
      a different encoder closing the gap, or a different LLM performing
      differently --- this design cannot distinguish between them. The
      agree/disagree confusion shared by B and C but absent from A
      (\autoref{sec:error-analysis}) tracks this same RoBERTa/LLM
      split, suggesting part of the ``strategy'' effect is really a
      backbone effect. Any claim involving A is scoped to this specific
      pairing, not a general LLM-vs-fine-tuning claim.
\item \textbf{Strategy C's seed variance was underestimated at $n=3$.} The
      original 3-seed standard deviation (0.0038) understated the true
      variance by roughly $10\times$ (0.037 across 8 seeds). Strategies A
      and B have not been checked with the same scrutiny.
\item \textbf{Held-out set, though expanded, is still drawn from the same
      corpus family as dev} (the Winning Arguments Corpus) --- not tested
      against a genuinely different conversational data distribution.
\item \textbf{Class imbalance in gold-300} (145 \texttt{disagree} vs.\ 64
      \texttt{agree}) is a candidate confound for the agree/disagree
      confusion finding that has not been ruled out.
\end{itemize}


% =====================================================================
% A-appendix.tex
% =====================================================================
\appendix
\chapter{Supplementary Tables}
\label{app:extra-matrices}

\begin{table}[htb]
\centering
\noindent\makebox[\textwidth]{%
\begin{tabular}{lrr}
\hline
\textbf{}&\textbf{Dev (gold-300)}&\textbf{Held-out (332)}\\ \hline
macro-F1& 0.5916& \textbf{0.6734}\\
Cohen's $\kappa$& 0.420& 0.558\\
95\% bootstrap CI& ---& [0.622, 0.720]\\
\hline
\end{tabular}
}
\caption{Strategy A: dev and held-out results}
\label{table:strategya-results}
\end{table}

\begin{table}[htb]
\centering
\noindent\makebox[\textwidth]{%
\begin{tabular}{lrrrr}
\hline
\textbf{True $\backslash$ Pred}&\textbf{agree}&\textbf{disagree}&\textbf{question}&\textbf{statement}\\ \hline
agree& 49& 3& 2& 29\\
disagree& 5& 48& 4& 26\\
question& 1& 12& 65& 5\\
statement& 9& 8& 6& 60\\
\hline
\end{tabular}
}
\caption{Strategy A: held-out confusion matrix (few-shot, 332 pairs)}
\label{table:strategya-heldout-cm}
\end{table}

\begin{table}[htb]
\centering
\noindent\makebox[\textwidth]{%
\begin{tabular}{lrr}
\hline
\textbf{}&\textbf{Dev (3 train seeds, mean)}&\textbf{Held-out (3 seeds, mean)}\\ \hline
macro-F1& 0.6375& \textbf{0.6686 $\pm$ 0.0012}\\
Cohen's $\kappa$& 0.4935& 0.5649\\
\hline
\end{tabular}
}
\caption{Strategy B: dev and held-out results (locked configuration)}
\label{table:strategyb-results}
\end{table}

\begin{table}[htb]
\centering
\noindent\makebox[\textwidth]{%
\begin{tabular}{lrrrr}
\hline
\textbf{True $\backslash$ Pred}&\textbf{agree}&\textbf{disagree}&\textbf{question}&\textbf{statement}\\ \hline
agree& 50& 12& 6& 15\\
disagree& 3& 51& 10& 19\\
question& 1& 3& 79& 0\\
statement& 17& 20& 2& 44\\
\hline
\end{tabular}
}
\caption{Strategy B: held-out confusion matrix (seed 123, representative)}
\label{table:strategyb-heldout-cm}
\end{table}

\begin{table}[htb]
\centering
\noindent\makebox[\textwidth]{%
\begin{tabular}{lrrrr}
\hline
\textbf{True $\backslash$ Pred}&\textbf{agree}&\textbf{disagree}&\textbf{question}&\textbf{statement}\\ \hline
agree& 47& 3& 1& 13\\
disagree& 19& 60& 17& 49\\
question& 1& 0& 43& 0\\
statement& 7& 10& 3& 27\\
\hline
\end{tabular}
}
\caption{Strategy B: dev confusion matrix (train\_seed 123, representative)}
\label{table:strategyb-dev-cm}
\end{table}

\begin{table}[htb]
\centering
\noindent\makebox[\textwidth]{%
\begin{tabular}{lr}
\hline
\textbf{Seeds}&\textbf{macro-F1 (mean $\pm$ std)}\\ \hline
Original 3 (42, 123, 2026)& 0.5309 $\pm$ 0.0038\\
All 8& \textbf{0.543 $\pm$ 0.037}\\
\hline
\end{tabular}
}
\caption{Strategy C: dev OOF variance, 3 vs.\ 8 seeds}
\label{table:strategyc-variance}
\end{table}

\begin{table}[htb]
\centering
\noindent\makebox[\textwidth]{%
\begin{tabular}{lrrr}
\hline
\textbf{}&\textbf{Dev OOF (3 seeds)}&\textbf{Dev OOF (8 seeds)}&\textbf{Held-out (3 seeds)}\\ \hline
macro-F1& 0.5309 $\pm$ 0.0038& 0.543 $\pm$ 0.037& \textbf{0.5859 $\pm$ 0.0699}\\
Cohen's $\kappa$& ---& ---& 0.483 $\pm$ 0.084\\
\hline
\end{tabular}
}
\caption{Strategy C: dev OOF vs.\ held-out results}
\label{table:strategyc-dev-vs-heldout}
\end{table}

\begin{table}[htb]
\centering
\noindent\makebox[\textwidth]{%
\begin{tabular}{lrrr}
\hline
\textbf{}&\textbf{A (few-shot)}&\textbf{B (silver 2500)}&\textbf{C (gold)}\\ \hline
Dev macro-F1& 0.5916& 0.6375& 0.531 (3 seeds) / 0.543 (8 seeds)\\
Held-out macro-F1 (332)& \textbf{0.673}& \textbf{0.669}& \textbf{0.586}\\
Held-out 95\% CI& [0.622, 0.720]& $\approx$[0.615, 0.719]& wide, seed-dependent\\
Old held-out (pilot80, $n{=}80$)& 0.699& 0.618& 0.569\\
\hline
\end{tabular}
}
\caption{Cross-strategy macro-F1 comparison}
\label{table:cross-strategy-comparison}
\end{table}

\begin{table}[htb]
\centering
\noindent\makebox[\textwidth]{%
\begin{tabular}{lrrr}
\hline
\textbf{Class}&\textbf{A}&\textbf{B}&\textbf{C}\\ \hline
agree& 0.667& 0.654& 0.580\\
disagree& 0.623& 0.590& 0.565\\
question& 0.812& 0.877& 0.837\\
statement& 0.591& 0.554& 0.362\\
\hline
\end{tabular}
}
\caption{Per-class F1 on held-out (332 pairs)}
\label{table:per-class-f1}
\end{table}

\begin{table}[htb]
\centering
\noindent\makebox[\textwidth]{%
\begin{tabular}{lrrrrr}
\hline
\textbf{Comparison}&\textbf{Observed $\Delta$}&\textbf{95\% CI of $\Delta$}&\textbf{Significant?}&\textbf{Power at $n{=}332$}&\textbf{$n$ needed (80\% power)}\\ \hline
A vs.\ B& 0.0041& [$-0.057$, 0.063]& No& 5\%& $\sim$34,900 pairs/class\\
A vs.\ C& 0.0465& [$-0.020$, 0.111]& No& 29\%& $\sim$330 pairs/class\\
B vs.\ C& 0.0424& [$-0.018$, 0.102]& No& 28\%& $\sim$350 pairs/class\\
\hline
\end{tabular}
}
\caption{Paired bootstrap power analysis ($n{=}332$, $n_{\text{boot}}{=}5000$)}
\label{table:power-analysis}
\end{table}

\begin{table}[htb]
\centering
\noindent\makebox[\textwidth]{%
\begin{tabular}{lrl}
\hline
\textbf{Strategy}&\textbf{Accuracy on 27 hard pairs}&\textbf{Own agree/disagree baseline}\\ \hline
A& 51.9\%& 49.8\% (no elevation)\\
B& 39.5\%& 55.7\% ($-16$ pts)\\
C& 0\% (by construction)& $\sim$47\%\\
\hline
\end{tabular}
}
\caption{Accuracy on the 27 ``hard'' agree/disagree pairs vs.\ each
strategy's own baseline}
\label{table:hard-pairs-accuracy}
\end{table}

\begin{table}[htb]
\centering
\noindent\makebox[\textwidth]{%
\begin{tabular}{lrrrr}
\hline
\textbf{True $\backslash$ Pred}&\textbf{agree}&\textbf{disagree}&\textbf{question}&\textbf{statement}\\ \hline
agree& 43& 3& 0& 17\\
disagree& 6& 60& 16& 62\\
question& 0& 5& 31& 7\\
statement& 2& 8& 3& 33\\
\hline
\end{tabular}
}
\caption{Strategy A: dev confusion matrix (few-shot)}
\label{table:strategya-dev-cm}
\end{table}

\begin{table}[htb]
\centering
\noindent\makebox[\textwidth]{%
\begin{tabular}{rrr}
\hline
\textbf{Silver size}&\textbf{Dev macro-F1 (mean $\pm$ std)}&\textbf{Runs}\\ \hline
500& 0.4155 $\pm$ 0.0865& 6\\
1000& 0.5814 $\pm$ 0.0189& 6\\
1500& 0.6143 $\pm$ 0.0275& 6\\
2000& 0.6184 $\pm$ 0.0187& 6\\
\textbf{2500}& \textbf{0.6220 $\pm$ 0.0274}& 6\\
5000& 0.6092 $\pm$ 0.0207& 6\\
8000& 0.5951 $\pm$ 0.0122& 6\\
10000& 0.5849 $\pm$ 0.0059& 3\\
\hline
\end{tabular}
}
\caption{Strategy B learning curve (few-shot-labeled silver pool, mean
$\pm$ std across sample\_seed $\times$ train\_seed runs)}
\label{table:strategyb-learning-curve}
\end{table}

\begin{table}[htb]
\centering
\noindent\makebox[\textwidth]{%
\begin{tabular}{lrrrr}
\hline
\textbf{Seed}&\textbf{agree row}&\textbf{disagree row}&\textbf{question row}&\textbf{statement row}\\ \hline
42& 37,29,11,6& 20,37,20,6& 3,0,78,2& 30,19,15,19\\
123& 59,18,2,4& 10,63,10,0& 1,1,79,2& 27,31,7,18\\
2026& 61,8,3,11& 15,51,12,5& 0,0,81,2& 27,25,4,27\\
\hline
\end{tabular}
}
\caption{Strategy C: held-out confusion matrices by seed}
\label{table:strategyc-heldout-cm-byseed}
\end{table}

\section{Strategy C: Full 8-Seed Dev OOF Confusion Matrices}

\begin{table}[htb]
\centering
\noindent\makebox[\textwidth]{%
\begin{tabular}{lrrrr}
\hline
\textbf{Seed}&\textbf{agree row [a,d,q,s]}&\textbf{disagree row}&\textbf{question row}&\textbf{statement row}\\ \hline
42& 24,30,0,10& 24,88,13,20& 1,5,35,3& 7,17,5,18\\
123& 34,20,1,9& 38,84,13,10& 1,6,35,2& 12,22,0,13\\
2026& 26,22,0,16& 19,95,16,15& 0,6,34,4& 5,26,1,15\\
7& 24,29,1,10& 26,98,12,9& 0,9,31,4& 12,17,1,17\\
99& 38,17,2,7& 27,93,13,12& 1,4,38,1& 9,21,3,14\\
777& 35,21,1,7& 21,95,13,16& 0,7,35,2& 7,23,0,17\\
1234& 35,16,1,12& 25,81,18,21& 0,2,41,1& 9,17,1,20\\
5555& 23,22,2,17& 40,69,16,20& 0,8,33,3& 11,16,1,19\\
\hline
\end{tabular}
}
\caption{Strategy C dev OOF confusion matrices, all 8 seeds (locked
configuration: warmup=20, epochs=8, weights on)}
\label{table:strategyc-8seed-devoof-cm}
\end{table}

In every one of the 8 seeds, the largest misclassification destination for
the \texttt{agree} row is \texttt{disagree}, and vice versa --- the basis
for the ``present in all 8 seeds without exception'' claim in
\autoref{sec:error-analysis}.

\section{Example Misclassified Pairs}

Representative concede-then-pivot examples referenced in
\autoref{sec:error-analysis} (true label $\to$ predicted label):

\begin{itemize}
\item \textbf{disagree $\to$ agree} (Strategy C): ``I understand that I'm
      simplifying...But my view is that right-wing ideas are inherently
      selfish...''
\item \textbf{disagree $\to$ agree} (Strategy C): ``Fair enough, you're
      entitled to your own opinion. Though I believe you're limiting
      yourself musically...''
\item \textbf{agree $\to$ disagree} (Strategy C): ``Based on that, which I
      agree with, you could also easily argue the case for Britney
      Spears.''
\end{itemize}
